% SPDX-License-Identifier: MIT-0
\documentclass[11pt]{article}

\usepackage[T1]{fontenc}
\usepackage[utf8]{inputenc}
\usepackage{lmodern}
\usepackage[a4paper,margin=2.4cm]{geometry}
\usepackage{amsmath,amssymb}
\usepackage{microtype}
\usepackage{booktabs}
\usepackage{longtable}
\usepackage[dvipsnames]{xcolor}
\usepackage{titlesec}
\usepackage[colorlinks=true,linkcolor=MidnightBlue,urlcolor=MidnightBlue,
            citecolor=MidnightBlue]{hyperref}

\hypersetup{%
  pdftitle={A manifest of the research reports in docs/reports},
  pdfauthor={Vladimir Reshetnikov},
  pdfsubject={Catalogue of 104 independent mathematical research packages}}

\definecolor{Ink}{HTML}{16324F}
\definecolor{Accent}{HTML}{4E5C3D}

\titleformat{\section}{\Large\bfseries\color{Ink}}{\thesection}{0.7em}{}
\titleformat{\subsection}{\large\bfseries\color{Accent}}{\thesubsection}{0.7em}{}

\newcounter{entrynum}
\setcounter{entrynum}{0}

\emergencystretch=1.6em

\def\UrlBreaks{\do\/\do\-\do\_\do\.}
\def\UrlBigBreaks{\do\/}

% \entry{title}{directory}{main document}{source archive}{description}
\newcommand{\entry}[5]{%
  \par\medskip\noindent
  \refstepcounter{entrynum}%
  \textbf{\theentrynum.\; #1}\par\nopagebreak
  \nopagebreak{\raggedright\small\color{Ink}%
    \nolinkurl{#2/}\textbf{\nolinkurl{#3}}\par}\nopagebreak
  \nopagebreak{\raggedright\footnotesize\color{gray}%
    from \texttt{#4}\par}\nopagebreak
  \smallskip\noindent #5\par}

\newcommand{\dir}[1]{\texttt{#1}}
\newcommand{\Q}{\mathbb{Q}}
\newcommand{\R}{\mathbb{R}}
\newcommand{\N}{\mathbb{N}}

\begin{document}

\begin{center}
  {\LARGE\bfseries\color{Ink} A manifest of \texttt{docs/reports}}\\[0.6em]
  {\large One hundred and four independent mathematical research packages,\\
   classified by subject}\\[1.0em]
  {\normalsize Vladimir Reshetnikov}\\[0.3em]
  {\normalsize 20 September 2026}
\end{center}

\vspace{0.8em}

\section{Scope and conventions}

This document catalogues the one hundred and four research packages stored
under \dir{docs/reports}.  The seven surreal-number reports that this catalogue
used to carry were moved out to a repository of their own, with their history,
and are no longer listed here.  They arrived as one hundred and twenty-five ZIP
archives, in seven deliveries, which have been unpacked and grouped into the
subject hierarchy described in Section~\ref{sec:tree}.  The archives themselves
are no longer kept in the repository.

One hundred and twenty-five archives yield one hundred and three reports for
two different reasons.  Twenty archives duplicated another: sixteen pairs, three
three-way clusters and one four-way cluster.  Two more were folded into a
neighbour not because they repeated it but because they shared a subject --- the
three reports on term counts in the derivatives of power towers are now one
article with a common framework and three parts.  That consolidation is
editorial, and the distinction matters: those three were put through the
duplicate sweep and cleared it, and their three theorems are three different
theorems.  The seventh delivery of eighteen contained no duplicate at all, which
is worth recording: a sweep that finds nothing is evidence about the collection,
not a sweep that failed.
Each pair has been merged into one report that proves the
shared theorem once and keeps everything both originals built on top of it.
Where the two reached a result by genuinely different routes, both proofs are
kept and marked as alternatives, with a sentence on what each one buys; the
merged shuffle report keeps three independent finite certifications of one
theorem, and the merged report on surreal option graphs kept two different
five-vertex counterexamples for the same reason before it moved out with the
other surreal-number reports.  A merged entry below names both of its
source archives and says what was doubled.

Not every pair aimed at one problem is a duplicate.  The two Gonshor
reports both prove his product-birthday bound, but on domains that are
incomparable in both directions, so neither theorem contains the other and they
are catalogued separately; the same holds for reports answering different numbered problems of
one paper, noted where it occurs.

Every package is self-contained and follows the same shape: a typeset article
(\texttt{*.pdf}) with its complete \LaTeX{} source, a \texttt{README.md}
summarising the result, usually a \texttt{code/} directory with exact-arithmetic
verification scripts, a \texttt{data/} or \texttt{results/} directory holding the
recorded output of those scripts, and often \texttt{STATUS.md},
\texttt{sources.md} or \texttt{notes/} files auditing the sources consulted and
the precise limits of the claim.

Two cautions apply uniformly and are therefore not repeated in the individual
entries.  First, these are AI-assisted research drafts: none has been
independently refereed, and none has been checked in a proof assistant.  Second,
the descriptions below record \emph{what each report sets out to settle and what
it claims to prove}.  Compiling this manifest involved reading each package's
statement of its problem and result; it did not involve checking any proof.
Where a report itself flags a limitation --- a prior disproof, an imported
theorem, a correction to the source, or an exception left open --- that is noted,
because it changes what the report is.

The naming of the directories is editorial.  Several distinct reports attack the
same published problem, and some arrived under colliding archive names, so each
directory has been given a name describing its own angle on the problem; the
merged directories are named for their combined scope, which is why none of them
still carries the name that once distinguished it from its partner.  The
original archive names are recorded with every entry.

\clearpage
\section{The directory tree}
\label{sec:tree}

\begin{center}
\begin{tabular}{@{}llr@{}}
\toprule
\textbf{Category} & \textbf{Subcategory} & \textbf{Reports}\\
\midrule
\dir{enumerative-combinatorics} & --- & 21\\
\midrule
\dir{ordinals-and-order-types} & \dir{ordinal-arithmetic} & 4\\
 & \dir{wqo-powersets-and-statures} & 4\\
 & \dir{friendly-order-types} & 3\\
 & \dir{games-on-ordinals} & 3\\
 & \dir{transfinite-words} & 2\\
 & (top level) & 4\\
\midrule
\dir{generating-functions-and-asymptotics} & \dir{oeis-sequence-asymptotics} & 8\\
 & \dir{binary-product-coefficient-counts} & 3\\
 & (top level) & 5\\
\midrule
\dir{log-concavity-and-unimodality} & --- & 12\\
\midrule
\dir{hankel-determinants} & \dir{catalan-and-ballot} & 4\\
 & \dir{somos-and-elliptic} & 3\\
 & \dir{growth-and-runs} & 2\\
 & (top level) & 1\\
\midrule
\dir{congruences-and-valuations} & \dir{supercongruences} & 4\\
 & \dir{iterated-series} & 2\\
 & (top level) & 3\\
\midrule
\dir{automata-and-formal-languages} & --- & 6\\
\midrule
\dir{tetration-and-digit-stabilization} & --- & 6\\
\midrule
\dir{graph-theory} & --- & 3\\
\midrule
\dir{quaternionic-analysis} & --- & 1\\
\midrule
\multicolumn{2}{l}{\textbf{Total}} & \textbf{104}\\
\bottomrule
\end{tabular}
\end{center}

\medskip
\noindent
More than a quarter of the collection consists of counterexamples and
corrections rather than proofs.  Ordinal Chomp, open-query games, compositional
tree series, the secant numbers, the decimal-length bound for tetration
stabilization, power-tower stabilization, iterated-power digit stabilization,
cluster variables, chromatic coefficients of cycles, products of $q$-integers,
domination roots, Hardinian array constants, the mixed-base counterexample,
recursive exponents, Gregory sign thresholds and binary substitution discrepancy
all refute, or replace with a proved statement, an assertion published or
catalogued elsewhere.  Several reports whose main outcome is a proof still
correct the statement they prove: the Somos-8 certificates, the shifted Catalan
Hankel polynomials and the Rueppel determinants record a sign error, a missing
one-unit shift and a numbering ambiguity in their sources, and the
independence-system thresholds report refines its target after finding that the
original conjecture had already been disproved by someone else.

\section{Ordinals and order types}

\subsection{Friendly order types}

Three reports on Isa Vialard's \emph{friendly order type} $f(P)$, the rank of
the tree of open-ended bad sequences in which every chosen point has an
incomparable partner in the current residual.  All three prove the same finite
formula; they differ in what they build on top of it, and in how they prove it
--- one by a forest-rank monovariant that survives moves deleting a whole upset,
one through a lemma bounding the cut maxima.  A fourth report proved the formula
by the first report's argument almost verbatim and has been merged into it.

\entry{Friendly Order Types Through Incomparability Components}
{ordinals-and-order-types/friendly-order-types/incomparability-components-and-infinite-tails}
{friendly\_order\_types.pdf}{friendly\_order\_type\_research (1).zip}
{The same component formula, together with an exact friendly rank for connected
infinite well-partial-orders: writing $o(C)=\lambda+n$ with $\lambda$ the limit
part and $T(C)$ the finite set of points with finite principal upsets.
Includes an exhaustive finite verification and a proof audit.}

\entry{Friendly Order Type Through Incomparability Components}
{ordinals-and-order-types/friendly-order-types/matroid-rank-and-ordinal-sums}
{article.pdf}{friendly\_order\_type\_research (3).zip}
{Identifies $|P|-c(P)$ with the graphic-matroid rank and the incidence-matrix
rank of the incomparability graph.  Adds a substitution formula, a finite
Cartesian-product formula, and additivity over arbitrary well-ordered ordinal
sums; evaluates every well-partial-order whose incomparability components are
finite, and gives countable counterexamples delimiting stronger infinite
extrapolations.}

\entry{Friendly Order Types of Finite Posets and Ordinal Products}
{ordinals-and-order-types/friendly-order-types/components-forests-and-ordinal-products}
{friendly\_order\_types.pdf}{friendly\_order\_types\_research.zip + friendly\_order\_type\_research.zip}
{Proves $f(P)=|P|-c(\mathrm{Inc}(P))$ for a finite poset $P$, where $c$ counts
connected components of the incomparability graph (isolated vertices included),
and classifies the witnesses: every maximum friendly subset is $P$ minus one
minimal root per component, counted by $\prod|\mathrm{Min}(C_i)|$.  With
Vialard's width-transfer theorem this gives ordinal width $\omega^{f(P)}$ for
the Dershowitz--Manna multiset ordering.  On the transfinite side it proves
$f(\alpha_1\times\cdots\times\alpha_r\times B)=\Theta$, the natural product,
except that it is the predecessor of $\Theta$ when every $\alpha_i$ is a
successor and $B$ has a greatest element, and classifies arbitrary finite
ordinal boxes.  Also: spanning-forest certificates, exact monotonicity under
added comparisons, a bivariate generating function, and explicit four-point
factors with equal $(o,h,w,f)$ whose products differ.  \emph{Merged; the
Cartesian-product formula is proved twice, once by corner deletion and once by
a chain-grid spanning-subgraph argument.}}

\subsection{Well-quasi-ordered powersets and statures}

Four reports on the open questions in Section~6 of Abriola, Halfon, Lopez,
Schmitz, Schnoebelen and Vialard, \emph{Measuring well-quasi-ordered finitary
powersets} (arXiv:2312.14587v2), and on the corresponding topological invariant.

\entry{Cofinal strata and ordinal absorption}
{ordinals-and-order-types/wqo-powersets-and-statures/cofinal-strata-of-finitary-powersets}
{article.pdf}{cofinal\_strata\_ordinal\_powersets.zip}
{For $S_\alpha(Q)$, a finite poset $Q$ with each point replaced by a chain
$\alpha=\omega^\rho$, proves
$o(K(S_\alpha(Q)))=\sum_k \omega^{\rho\otimes k}c_k(Q)$ and
$h(K(S_\alpha(Q)))=\alpha\cdot\mathrm{height}(M(Q))$, where $c_k$ counts size-$k$
maximal antichains not dominated by a larger one.  The maximal-order-type
theorem allows nonuniform pure well-partial-order fibers and is proved by a
finite cofinal-stratification argument.}

\entry{Intrinsic ordinal approximations for finitary powersets}
{ordinals-and-order-types/wqo-powersets-and-statures/intrinsic-ordinal-approximations}
{article.pdf}{intrinsic\_ordinal\_approximations.zip}
{Studies the intrinsic invariant
$a_o(A)=\sup\{o(B)+1: \mathrm{Idl}(B)\hookrightarrow A\}$.  Proves the universal
bound $a_o(A)\le h(P_f(A))$ with equality exactly when every finitely generated
downset of $A$ is an algebraic dcpo, gives the closed formula
$h(P_f(\prod_i\alpha_i))=\sup(\bigotimes_i(1+a_i)+1)$, and shows the universal
identity fails on $\lambda+F$.}

\entry{Ordinal Order Maps and Finitary Powersets of Ordinal--Finite Grids}
{ordinals-and-order-types/wqo-powersets-and-statures/ordinal-order-maps-and-grids}
{ordinal\_order\_maps.pdf}{ordinal\_order\_maps\_research.zip}
{Exact maximal order types and heights for finite Hoare powersets of
$\beta\times Q$ with $Q$ a finite poset: for
$\beta=\omega^{\gamma_0}+\cdots+\omega^{\gamma_{m-1}}$,
$o(\mathrm{Mon}(Q,\beta))=\bigoplus_{s\,\text{isotone}}\omega^{\bigoplus_q\gamma_{s(q)}}$.}

\entry{Sharp ordinal bounds for Hoare powerspaces}
{ordinals-and-order-types/wqo-powersets-and-statures/hoare-powerspace-statures}
{article.pdf}{hoare\_statures\_research.zip}
{For all Noetherian spaces, $1+\|X\|\le\|HX\|\le 2^{\|X\|}$, with witnesses
attaining both endpoints at every ordinal and $T_1$ upper witnesses; the older
$\omega^\alpha$ bound is attainable at infinite input stature $\alpha$ exactly
when $\omega\cdot\alpha=\alpha$; the complete spectrum at $\omega+n$ is
$\omega\cdot j(U)+(j(Q)-j(U))$.  Extends the known well-quasi-order results of
Abriola et al.\ to arbitrary Noetherian spaces.}

\subsection{Transfinite words}

Two reports on Harry Altman, \emph{Bounding finite-image sequences of length
$\omega^k$} (arXiv:2409.03199v2), Example~3.15: is the bound
$\omega^{\omega^4}$ sharp for words of length ${<}\,\omega^2$ over the
two-element antichain?  A third report matched the first lemma for lemma and has
been merged into it.

\entry{Exact maximal order types of finite-alphabet transfinite words}
{ordinals-and-order-types/transfinite-words/finite-alphabet-transfinite-words}
{article.pdf}{finite\_alphabet\_transfinite\_words.zip}
{Proves $o(W_k(P))=\omega^{\omega^{N_k(P)-1}}$ for words of ordinal length
${<}\,\omega^k$ over a finite poset $P$ under subsequence embedding, where
$N_k$ counts indecomposable classes of a canonical finite atom poset
$I_k(P)$, with $N_1(P)=|P|$ and $N_{k+1}(P)=|P|+J(I_k(P))-1$.  For $k=2$ the
binary antichain gives $\omega^{\omega^4}$ --- the tightness Altman leaves
open --- and the binary chain gives $\omega^{\omega^3}$.}

\entry{Guarded periodic blocks and exact maximal order types below omega squared}
{ordinals-and-order-types/transfinite-words/order-types-below-omega-squared}
{article.pdf}{omega2\_order\_types\_proposed\_solution.zip + ordinal\_words\_research.zip}
{Proves $o(s_{\omega^2}(P))=\omega^{\omega^{n+j(P)-2}}$ for $n=|P|\ge2$, with
$j(P)$ the number of downsets including the empty one; the empty and singleton
alphabets are genuine exceptions.  The binary antichain gives $\omega^{\omega^4}$
--- the tightness Altman leaves open --- and the binary chain $\omega^{\omega^3}$.
Adds an extremality corollary with both equality characterisations (among
$n$-element posets only the chain minimises and only the antichain maximises),
the $j(P)$ toolkit, the arbitrary-support-family classification, and an
exact-length stratum calculus showing the supremum over lengths is only
$\omega^{\omega^n}$, strictly below the answer --- so the theorem cannot be
obtained stratum by stratum.  \emph{Merged; the product amplification is done
twice, once via de Jongh--Parikh and once by an elementary hand proof that
removes the import.}}

\subsection{Ordinal arithmetic}

Four reports on ordinal and infinitary arithmetic, two of them independent
proposed solutions to Problem~6.2 of Paolo Lipparini,
\emph{Monotone infinitary operations on ordinals} (arXiv:2505.00424v2): define
and study the least weakly monotone, $e$-special strictly increasing operation
on countable sequences of ordinals.

\entry{An Explicit Formula for the Minimal $e$-Monotone Infinitary Ordinal Sum}
{ordinals-and-order-types/ordinal-arithmetic/lipparini-minimal-infinitary-sum}
{article/explicit\_ordinal\_sum.pdf}{lipparini\_problem\_6\_2\_proposed\_solution.zip}
{An explicit formula for the least admissible operation, proved by an arithmetic
admissibility upper bound and a rank-theoretic minimality lower bound, using an
explicit countable-multiset profile order, its rank, absorption and correction
rules, bounded-poset heights and finite-cap ranks.  Imports Lipparini's
formula and minimality theorem for the weaker operation $S$.}

\entry{An Explicit Formula for Lipparini's Minimal Infinitary Ordinal Operation}
{ordinals-and-order-types/ordinal-arithmetic/lipparini-minimal-operation-formula}
{article.pdf}{ordinal\_operation\_research.zip}
{A three-case finite ordinal-arithmetic formula with proofs of admissibility and
pointwise minimality.  Relative to Lipparini's already evaluated $S$, the
correction is zero, a finite increment, or one \emph{natural} summand $\omega$
--- ordinary addition of $\omega$ is not a substitute.}

\entry{Jacobsthal Products All at Once}
{ordinals-and-order-types/ordinal-arithmetic/jacobsthal-products-all-at-once}
{article.pdf}{Jacobsthal\_Products\_All\_at\_Once.zip}
{Question~2.16 of Altman, \emph{Intermediate arithmetic operations on ordinal
numbers} (MLQ 63 (2017) 228--242), asks for an all-at-once order-theoretic model
of finite and infinitary Jacobsthal products, rather than one built from
parenthesized binary multiplication or a limit.  Gives a constructive
affirmative answer: a Cantor-normal-form comparison rule on actual tuples or
finite-support functions, invariant under order isomorphisms and extending the
componentwise order, whose definition never evaluates a Jacobsthal product.
Adds a restricted maximality characterisation and a finite-tail reduction.
The associativity identity and the binary normal-form formula are classical and
are not claimed here; the construction is not claimed to be preserved by
arbitrary order embeddings.}

\entry{Pairwise Coherence Does Not Guarantee Completion}
{ordinals-and-order-types/ordinal-arithmetic/pairwise-coherence-of-infinitary-products}
{ordinal\_product\_coherence.pdf}{ordinal\_product\_coherence.zip}
{Answers the pairwise-coherence sufficiency question of Lipparini,
\emph{Noncommutative infinitary semigroups} (arXiv:2605.28413v1, Remark~7.1(b)),
negatively.  On the two-element group $C_2$, retaining every finite product and
declaring every infinite all-$e$ product to be $e$ gives a partial system
satisfying full pairwise regrouping, even for arbitrary partitions, yet with no
injective extension defining all finite and $\omega$-indexed products; two
elements is minimal.  Adds: for any monoid $M$ the canonical partial domain
$P(M)$ is coherent, and for directly finite $M$ it completes exactly when the
unit group is trivial; and for every $r$ a system on $3r+4$ elements coherent
through nesting depth $r$ but conflicting at depth $r+1$, so no fixed finite
depth of the coherence test characterises completability.}

\subsection{Games on ordinals}

\entry{An ordinal Chomp transition at 2}
{ordinals-and-order-types/games-on-ordinals/ordinal-chomp-transition-at-two}
{ordinal\_chomp\_counterexample.pdf}{ordinal\_chomp\_transition\_counterexample.zip + ordinal\_chomp\_counterexample.zip}
{For a numerical semigroup $S$ ordered by $x\le_S y$ iff $y-x\in S$, and the
ordinal extension ordered by ordinary ordinal subtraction (so $S^2$ is the
\emph{lexicographic}, not coordinatewise, square), Chomp on $S$ is a
second-player win while Chomp on $S^\sigma$ is a first-player win for every
$\sigma\ge2$, and $\omega\cdot4+4$ is the \emph{least} winning opening.  Hence
$\mathrm{ch}(\langle4,5,6\rangle)=\mathrm{ch}(\langle4,6,9\rangle)=2$, a negative
answer to Question~4.6 of Rivero Herrera (arXiv:2504.07317v3) and a refutation of
Conjecture~4.1 of v1, already for three small natural-number generators.  The
proof reduces the infinite game to a finite spectral statement via a
finite-window recurrence and an exact finite-state certificate, then transfers
it by an explicit response procedure.  \emph{Merged; the shared
$\langle4,6,9\rangle$ certificates were verified to hold the same 253 rows
entry for entry, and the two reports' poisoned and unpoisoned conventions are
reconciled by an explicit dictionary.}}

\entry{Every ordinal as a point-separating game value}
{ordinals-and-order-types/games-on-ordinals/point-separating-game-values}
{ordinal\_separation.pdf}{ordinal\_separation\_research.zip}
{For every ordinal $\lambda$, orders $\{0,1\}^\lambda$ lexicographically and
gives it the topology of all downward-closed sets, and the refinement generated
by downward-closed sets minus finite sets.  The point-separating game value is
exactly $\lambda$ for both: the first space is compact $T_0$, the second compact
$T_1$, and for infinite $\lambda$ neither is Hausdorff.  Taking
$\lambda=\omega^2+1$ supplies a compact $T_1$ witness of cardinality continuum
for Problem~1.2 of Chiozini, Csern\'ak and Soukup (arXiv:2510.05754v3), with no
CH, forcing or large-cardinal hypothesis.  Hausdorff witnesses are not claimed.
This is a different problem of that paper from the one the open-query games
report answers.}

\entry{Open-query games on ordinal spaces}
{ordinals-and-order-types/games-on-ordinals/open-query-membership-games}
{article.pdf}{ordinal\_membership\_games.zip}
{A negative answer to Problem~1.3 of Chiozini, Csern\'ak and Soukup
(arXiv:2510.05754v3): two copies of $[0,\omega]$ each have set-membership
number~1, but their topological sum has value~2.  The general theorem computes
the invariant for every nonempty Hausdorff space of finite Cantor--Bendixson
height~$h$ as $\lceil\log_2(h+1)\rceil$ when the last nonempty derivative is a
singleton and $\lceil\log_2(h+2)\rceil$ otherwise.}

\subsection{Order-theoretic invariants without a subcategory}

Two reports whose objects sit outside every subcategory above: the minor order
on finite graphs, and a site of filters on $\omega$.

\entry{Ordinal complexity of graph minors}
{ordinals-and-order-types/graph-minor-ordinals}
{graph\_minor\_ordinals.pdf}{graph\_minor\_ordinals.zip}
{For finite simple graphs of maximum degree at most two under minor
containment, computes maximal order type, ordinal antichain-tree width and
descending-tree height: $o(D)=w(D)=\omega^{\omega\cdot2}$ and $h(D)=\omega$,
with exact values $\omega^{b+k}$ and $\omega^{b+k-1}$ for the subfamilies
bounded by path size $b$ and cyclic-component count $k$.  The general theorem is
that a wqo $C$ of connected graphs with $o(C)=\alpha<\varepsilon_0$ has finite
disjoint-union closure of type $\omega^\alpha$, proved by a component-splitting
argument, since minor containment is not multiset embedding.  A corollary is
that continuity fails for increasing unions of minor-closed classes.  The
report is explicit that it does not compute the maximal order type of the minor
order on all finite graphs and settles no named conjecture.}

\entry{A two-valued obstruction to gluing Fr\'echet powers}
{ordinals-and-order-types/frechet-two-valued-obstruction}
{article.pdf}{frechet\_two\_valued\_obstruction.zip}
{For the presheaf $F\mapsto A^\omega/{\sim_F}$ on proper filters on $\omega$
containing the cofinite filter, with the dense Grothendieck topology, proves in
ZF that it is a sheaf if and only if $|A|\le1$.  The counterexample is an
explicit matching family of constant sections using two values, on a covering
sieve with countably many generators each of countable base; the distinction
between the number of sieve generators and the size of the full sieve is
essential to it.  The question is Massas, \emph{A Semi-Constructive Approach to
the Hyperreal Line} (2023), p.~517 after Fact~5.10, repeated in the 2024
dissertation \emph{Duality and Infinity}, p.~249.}

\entry{A sharp cover bound for congruences of finite lattices}
{ordinals-and-order-types/lattice-congruence-cover-bound}
{sharp\_cover\_bound.pdf}{lattice\_cover\_bound\_research.zip}
{If a finite $n$-element lattice has an element with at least $k\ge3$ upper
covers then $|\mathrm{Con}(L)|\le2^{\,n-k-1}$, with equality exactly for the
glued sums $C_s\oplus M_k\oplus C_t$.  This replaces the multicolour Ramsey
threshold $R_{k-1}(k)$ of the source's Lemma~4.4 (arXiv:2603.11454) by $k$
itself.  Adds a sharp nonextremal gap, a fan-label ideal refinement, and a
bound on the size of Cz\'edli's skeleton in terms of congruence density.  All
congruence machinery, including Freese's maximum, is proved from scratch.  The
target is an informal suggestion in the source, not a numbered conjecture.}

\entry{A non-Baire translation-invariant ideal}
{ordinals-and-order-types/non-baire-translation-invariant-ideal}
{article.pdf}{non\_baire\_ideal\_research.zip}
{An affirmative construction for Question~1 of Filipów and Tryba,
\emph{Densities for sets of natural numbers vanishing on a given family}
(J.\ Number Theory 211 (2020) 371--382; arXiv:2309.00982, p.~7): a
translation-invariant ideal that is not Baire and carries a perfect
translation-almost-disjoint family.}

\section{Automata and formal languages}

Five reports on state complexity, synchronization and language hierarchies.
The first two are each a merge: the shuffle bound arrived four times and the
DFAO reversal bound twice.

\entry{The six-state shuffle bound}
{automata-and-formal-languages/shuffle-six-state-bound}
{shuffle\_six\_state.pdf}{shuffle\_six\_state\_research (1).zip + shuffle\_six\_state\_research.zip + Shuffle\_Six\_State\_Theorem.zip + six\_state\_shuffle\_research.zip}
{For complete DFAs with $m,n\ge2$ and $\min(m,n)\le6$, the worst-case
deterministic state complexity of the shuffle is
$F(m,n)=2^{mn-1}+2^{(m-1)(n-1)}(2^{m-1}-1)(2^{n-1}-1)$, over a finite alphabet
allowed to depend on $m$ and $n$; the unrestricted conjecture stays open.  The
route is a reachability theorem in the universal shuffle automaton --- a subset
is reachable exactly when it meets the initial row and the initial column ---
proved by induction on $(m+n,|S|)$ down to a finite certificate obligation.
\emph{Merged from four reports.  Two shipped the same 10{,}642 witnesses; a
third supplies an independent database of 10{,}619 witnesses sharing no witness
with them, so the finite step is certified twice over; a fourth adds a fifth
analytic reduction that cuts the obligation to 453 certificates and replaces
the solver search with a deterministic enumeration.  All three routes are
kept.}}

\entry{A coloring obstruction to maximal binary DFAO reversal}
{automata-and-formal-languages/dfao-reversal-coloring-obstruction}
{article.pdf}{dfao\_reversal\_research.zip + binary\_dfao\_reversal\_research.zip}
{For transformations $a,b$ on an $n$-element set and an output map to $k$
labels, proves $|\tau\langle a,b\rangle|\le k^n-k!+g(k)<k^n$ for $3\le k\le n$,
with $g(k)$ the largest order of a permutation of $k$ symbols --- the
binary-alphabet nonattainment question of Davies, arXiv:1705.07150v2,
Section~5.  The mechanism is a collision-orbit graph, a rightmost singular
factorization, and a cyclic relabeling bound.  Adds an abelian-permutation
extension, a structural universal bound matching the published lower
construction on all 372 pairs with $7\le n\le30$, and polynomial-time
missing-coloring certificates.  \emph{Merged; the two originals ran the same
three-case argument, and both a qualitative short proof and the counting proof
are kept.}}

\entry{Synchronizing D\.zyga's automata by a cyclic retract}
{automata-and-formal-languages/dzyga-synchronizing-retract}
{article.pdf}{Dzyga\_Automata\_Cyclic\_Retract\_Research.zip}
{For the Figure-3-complete automaton $A_k$ of Conjecture~8 in Szyku\l{}a's 2026
survey, exhibits an explicit word $R_k=p(ec^{r-1})^{r-2}e$ with $r=k+4$ that
sends every state to $q_2$ for every $k\ge1$, of expanded length
$8k^2+54k+92=\tfrac89n^2+\tfrac{22}3n+16$ for $n=3k+6$.  The proof is a cyclic
retraction with two projected reflections whose discrepancy supplies one
adjacent merger, and does not depend on the finite computations; shortening the
reflected tail by one state gives a nonsynchronizing automaton with a dihedral
permutation quotient.  This settles the synchronization clause of the stated
model, not the numerical threshold conjecture, and the universal lower bound
$\tau(q_0)\ge4k+8$ is not proved.}

\entry{Gaps in bounded-shuffle hierarchies}
{automata-and-formal-languages/insertion-degree-spectra}
{article.pdf}{shuffle\_spectra\_research.zip}
{A counterexample to the insertion-degree No-Gap Conjecture of Hughes,
arXiv:2608.27755v1, Section~11, followed by positive results: every finite set
of positive integers containing~1 is the insertion-degree spectrum of two
nonempty finite languages over a suitable alphabet, and any prescribed finite
list of degrees above~1 is realizable with every other output of degree~1.  A
ternary pair has spectrum $\{1,3\}$ with sole difficult output \texttt{abcba},
whose length~5 is the minimum possible for a gap witness; a binary pair has
spectrum $\{1,2,4\}$ while all 18 singleton source-pair spectra are initial
intervals, so source-pair ambiguity can erase a degree.  The separate
iteration-depth no-gap conjecture is proved for arbitrary languages.}

\entry{A cut-sensitive attack on the first open MIX grammar case}
{automata-and-formal-languages/mix-grammar-cut-sensitive}
{mix\_grammar\_report.pdf}{mix\_grammar\_research.zip}
{The selected question --- whether $L_3=\mathrm{MIX}$ for the canonical tuple
grammars of Yuyama, \emph{On the Kanazawa--Salvati Conjecture}
(arXiv:2609.13871v1) --- is \emph{not} settled here, and the report says so at
the front.  What it does contain, with complete written proofs, is a
cut-sensitive membership recurrence, a broad three-component positive family,
an infinite family in $L_3\setminus L_2$, and an exhaustive computation through
length~21.}

\entry{Natural densities of the Tribonacci word's additive complexities}
{automata-and-formal-languages/tribonacci-additive-complexity}
{article.pdf}{tribonacci\_additive\_density\_research.zip}
{Answers Remark~19 of Popoli, Shallit and Stipulanti (arXiv:2410.02409).  With
$\beta$ the Tribonacci constant, the additive complexity of the Tribonacci word
takes values in $\{3,4,5\}$ with natural densities
$d_3=(21\beta^2-38\beta+5)/22$, $d_4=(-55\beta^2+72\beta+67)/22$ and
$d_5=(17\beta^2-17\beta-25)/11$, all positive and summing to~1, with a
power-saving error $O(N^\theta)$, $\theta=\log(7/5)/\log\beta$, at
\emph{every} integer cutoff rather than only along Tribonacci cutoffs.  The
76-state output automaton is reconstructed independently rather than copied,
and certified by evaluating its degree-76 characteristic polynomial at the
adjacency matrix.}

\section{Hankel determinants and continued fractions}

\subsection{Somos sequences and elliptic curves}

\entry{A proof of Barry's Hankel--Somos-6 conjecture}
{hankel-determinants/somos-and-elliptic/barry-somos-6-hankel-recurrence}
{article.pdf}{barry\_somos6\_research.zip}
{Conjecture~7 of Paul Barry, arXiv:2211.12637: for the Hankel determinants
$H_N$ of a four-parameter Jacobi-type generating function,
$H_NH_{N-6}=\alpha H_{N-1}H_{N-5}+\gamma H_{N-3}^2$ with
$\alpha=t^2(r+2)^2$ and
$\gamma=t^3((s+t)(r+s+t+2)^2+t(r+2)^3)$.  Also proves a homogeneous
four-parameter form and a seven-parameter polynomial Hankel theorem making
exceptional specialisations rigorous; the verifier never divides by a
determinant.}

\entry{Finite Hankel Certificates for Three Somos-8 Conjectures}
{hankel-determinants/somos-and-elliptic/somos-8-hankel-certificates}
{article.pdf}{somos\_hankel\_proofs.zip}
{Conjectures~11, 12 and~13 of the same Barry paper, by importing Hone's
genus-two Hankel theorem (arXiv:1907.05204v3, Thm.~5.4) and supplying a
binomial translation, a polynomial-minor specialisation valid on the degenerate
locus, and a rank-bound lemma that turns four symbolic rows into an all-index
identity.  Identifies a sign error in the source's Catalan expression for
Conjecture~13.}

\entry{Elliptic addition and Jacobi continued fractions}
{hankel-determinants/somos-and-elliptic/elliptic-addition-and-jacobi-fractions}
{article.pdf}{elliptic\_continued\_fraction\_research.zip}
{A consistently indexed, non-torsion version of Barry's 2023 Conjecture~4, with
the coordinate formulas of Conjecture~3: every continued-fraction tail is given
explicitly and removing one layer is shown to be elliptic translation.
Hypotheses: characteristic zero, $Y^2+aXY+bY=X^3+cX^2+dX$ nonsingular,
$b\ne 0$, and $P=(0,0)$ of infinite order; order-three and order-four examples
show the regular fraction genuinely needs the last one.}

\subsection{Catalan and ballot moments}

\entry{Shifted Catalan Hankel Polynomials}
{hankel-determinants/catalan-and-ballot/shifted-catalan-hankel-polynomials}
{article.pdf}{catalan\_hankel\_conjecture\_solution.zip}
{For $D_N^{(m)}(a,b)=\det(aC_{i+j+m}+bC_{i+j+m+1})$, proves the
denominator-exponent assertion on page~4 of Barry, arXiv:2011.10827v1, and
shows the separately printed Conjecture~2 needs an explicit one-unit shift
correction.  The minimal characteristic polynomial is
$(X^2-(a+2b)X+b^2)^{m(m-1)/2+1}$, so the minimal order is $m(m-1)+2$; all
exceptional reduced denominators, the exact coefficient formula, numerator
degree and reciprocal symmetry are supplied.}

\entry{Product formulas for ballot-polynomial Hankel determinants}
{hankel-determinants/catalan-and-ballot/ballot-polynomial-hankel-determinants}
{article.pdf}{Cigler\_Hankel\_Conjectures\_13\_15.zip}
{Conjectures 13, 14 and 15 in Section~5 of Johann Cigler, arXiv:2111.14492v3.
An explicit coefficient product gives the sign and exact support of the
normalised shift-$k$ Hankel determinant of the ballot moments and positivity of
every residual coefficient for $n\ge1$; the same products give the conjectured
generating-function denominators, exact numerator degrees and reciprocal
numerator identities.  Adds strict log-concavity of the residual coefficients
and their limiting binomial profile.}

\entry{Rectangular Schur polynomials and Cigler's Conjecture 8}
{hankel-determinants/catalan-and-ballot/cigler-conjecture-8-schur}
{article.pdf}{cigler\_conjecture8\_research.zip}
{A complete proof of Conjecture~8 in Cigler, arXiv:2111.14492v3, Section~4
(printed pp.~16--17), for the Hankel determinants of
$b_m(t)=\sum_r\binom{\lfloor m/2\rfloor}{r}\binom{\lceil m/2\rceil}{r}t^r$,
including the numerator-polynomial degree and reciprocity assertions.  A
different conjecture of that paper from the ballot-polynomial entry
(Conjectures 13--15) and from the parity entry (Conjecture~16).}

\entry{A parity factorization for Cigler's Conjecture 16}
{hankel-determinants/catalan-and-ballot/cigler-conjecture-16-parity}
{article.pdf}{cigler\_conjecture16\_research.zip}
{A complete proof of Conjecture~16 in Cigler, arXiv:2111.14492v3, Section~6
(pp.~21--22), on the Hankel polynomials, together with a sharp stabilization
statement: the report also determines the first coefficient at which the
conjectured stabilization fails.  Its own README explains why this is a
different problem from the other Cigler entries in this collection.}

\subsection{Arithmetic Hankel determinants}

The numerator formula recorded as a conjecture in OEIS A122251, with A122249 as
its binary special case.  Two reports attacked it and proved it the same way,
sharing even their worked counterexamples; they are merged here into one.  The
merged report sits at the category's top level.

\entry{A Proof of an OEIS Hankel-Numerator Conjecture}
{hankel-determinants/a122251-numerators-and-denominators}
{report.pdf}{hankel\_numerator\_research.zip + hankel\_oeis\_conjecture\_proof.zip}
{For $H_n(m,a)=\det(1/(a+m(i+j)))_{0\le i,j\le n}$, the reduced numerator is
$\prod_{p\mid m}p^{\,n(n+1)v_p(m)+2\sum_{k\le n}v_p(k!)}$: independent of the
coprime offset $a$, a perfect square, and supported on the primes dividing $m$.
Taking $m=p$ and $a=1$ settles the OEIS A122251 conjecture, with A122249 as the
binary case.  Also gives a closed form for the reduced denominator and all its
local valuations, a sharp gcd theorem over the allowed offsets realised at just
two explicitly constructed shifts, the extension to unequal row and column steps
$\det(1/(a+m_1i+m_2j))$, recurrences, digit-sum asymptotics, and a Lambert-series
generating function with a natural boundary.  \emph{Merged; the determinant is
evaluated three independent ways, and a second proof of the numerator formula
runs through an explicit inverse matrix and MacMahon's boxed plane partitions.}}

\subsection{Growth laws and run statistics}

\entry{A proof of the Ap\'ery Hankel growth conjecture}
{hankel-determinants/growth-and-runs/apery-hankel-determinant-growth}
{apery\_hankel\_growth.pdf}{apery\_hankel\_research.zip}
{For $A_k=\sum_j\binom{k}{j}^2\binom{k+j}{j}^2$ and $D_n=\det(A_{i+j})$, proves
$\log D_n=n(n+1)\log\frac{17+12\sqrt2}{4}+O(n)$ --- the limit conjectured in the
September 2026 revision of OEIS A228143 --- using Edgar's moment-density
theorem.  Also an explicit law for $\det(A_{m(i+j)+r_n})$ with $r_n/n\to\rho$,
and non-D-finiteness of the formal generating series of $D_n$.}

\entry{Binary Runs in Rueppel Hankel Determinants}
{hankel-determinants/growth-and-runs/rueppel-binary-run-determinants}
{rueppel\_hankel.pdf}{rueppel\_hankel\_research.zip}
{For the Rueppel series $r(x)=\sum_{j\ge0}x^{2^j-1}$, the polynomial identity
$H_n(tx+1/r(x^2))=(-1)^{1+\varepsilon+\lfloor R(n)/2\rfloor}
(\varepsilon+(R(n)-\varepsilon)t^2)$, where $R(n)$ counts runs of equal bits in
the binary expansion of $n$ and $\varepsilon$ is the low bit of
the Gray code $n\oplus\lfloor n/2\rfloor$.  Settles statements 8, 10, 13 and 23 of Barry's 2021 paper
(not the identically numbered ones elsewhere), with a correction.}

\section{Congruences and $p$-adic valuations}

\subsection{Supercongruences}

\entry{Congruences for OEIS A348410 at every prime}
{congruences-and-valuations/supercongruences/a348410-all-prime-congruences}
{article.pdf}{a348410\_supercongruence.zip + A348410\_Supercongruence\_Research.zip}
{For $A_{\alpha,\beta}(n)=[x^n](1-x)^{-\alpha n}(1+x)^{-\beta n}$ proves the
odd-prime cubic supercongruence $v_p(A(Np)-A(N))\ge 3v_p(N)-\varepsilon_p$ with
$\varepsilon_3=1$, $\varepsilon_p=0$ for $p\ge5$, and no coprimality condition;
$(\alpha,\beta)=(2,1)$ is Bala's A348410 conjecture.  Separately proves an exact
binary leading-defect formula, $A(m2^r)-A(m2^{r-1})\equiv
2^{3r-3}\alpha\beta\binom{(\alpha+\beta+1)m-1}{m-1}$ mod $2^{3r-2}$, hence the
all-prime bound $3r-v_p(12)$ and sharp denominators $12$ and $24$.  An appendix
audits arXiv:2104.10754v1, whose framing conclusion and weighted harmonic
assertion fail on its own displayed examples.  \emph{Merged; the odd-prime and
binary halves run on different mechanisms --- a doubling automorphism, and
reflection with formal integration by parts --- and both are kept.}
\par\smallskip\noindent\textbf{Not a duplicate of
\emph{signed-binomial-supercongruences}, and checked.}  That report's
$a_{r,s,t}(n)=[x^{tn}](1+x)^{rn}(1-x)^{sn}$ specializes at $t=1$,
$(r,s)=(-\beta,-\alpha)$ to the coefficient studied here, and the two
odd-prime statements then agree.  Neither contains the other.  That one is
general in the extraction step $t$, which is fixed at $1$ here; this one adds
the prime $2$ with an exact defect and its sharpness, every prime at once,
minimal common denominators, the primitive-cyclic-word interpretation, the
Lambert/Euler product, the algebraic generating functions, a
spacing-three counterexample, and a sufficient criterion beyond the family.
The proofs are different arguments: that one imports Jacobsthal's congruence,
this one uses none and pairs a quadratic-tail estimate with a coefficient
estimate.
}

\entry{Signed binomial supercongruences and OEIS A352373}
{congruences-and-valuations/supercongruences/a352373-signed-binomial-supercongruences}
{signed\_binomial\_supercongruence.pdf}{OEIS\_A352373\_supercongruence.zip}
{For $a(n)=[x^{tn}](1+x)^{rn}(1-x)^{sn}$ proves
$a(mp^h)\equiv a(mp^{h-1})\bmod p^{3h}$ for $p\ge5$ and
$a(m3^h)\equiv a(m3^{h-1})\bmod 3^{3h-1}$, with $m$ not required coprime to
$p$.  Covers A352373, A348410, A351856 and A351857 as specialisations; the only
non-elementary import is the classical Jacobsthal binomial congruence.
\par\smallskip\noindent\textbf{Not a duplicate of
\emph{all-prime-congruences}, and checked.}  Setting $t=1$ and
$(r,s)=(-\beta,-\alpha)$ turns the theorem below into that report's
odd-prime tower, and the moduli agree after reindexing the multiplier.
Neither contains the other: this one is general in the extraction step $t$,
that one is fixed at $t=1$ but proves the prime~$2$ case, the all-prime
bound, the sharp denominators and much else.  The two proofs are genuinely
different --- this one imports Jacobsthal's congruence, that one does not use
it at all.
}

\entry{Apéry numbers under exponentiation and series reversion}
{congruences-and-valuations/supercongruences/a267220-apery-transform-supercongruences}
{article.pdf}{Apery\_OEIS\_Supercongruences.zip, apery\_oeis\_conjectures(1).zip}
{Proves both supercongruence conjectures contributed by Peter Bala to OEIS
A267220 on 17 October 2024, for every integer parameter including negative and
zero.  With $A(x)=\exp\sum_{n\ge1}a_nx^n/n$ for the Apéry numbers $a_n$,
$F(x)=x/\mathrm{Rev}(xA(x))$, $u_m(n)=[x^n]A(x)^{mn}$ and
$v_m(n)=[x^n]F(x)^{mn}$: for odd $p$ and $R=\nu_p(N)$ both families satisfy
$\nu_p(u_m(N)-u_m(N/p))\ge2R+\nu_p(m)$, the unsigned statement for $v_m$ holds
at \emph{every} prime including $2$, and a \emph{signed} refinement
$(-1)^{mN}u_m(N)\equiv(-1)^{mN/p}u_m(N/p)$ recovers full precision at $p=2$,
where the unsigned bound genuinely loses one power.  Merged from two reports
that share a seed lemma and a Möbius/dilogarithm layer but diverge at the crux;
both routes are kept, and the merged text is careful that the signed and
unsigned $2$-adic statements are different statements, not a sharper and a
weaker version of one.}

\entry{Four Motzkin congruences}
{congruences-and-valuations/supercongruences/a001006-motzkin-congruences}
{article.pdf}{motzkin\_congruences.zip}
{Proves all four conjectures in Peter Bala's comment of 10 February 2022 on OEIS
A001006, with sharp ranges.  The engine is a prime-power block-transfer theorem:
for odd $p$ and $q=p^{k-s+1}$, $h=mp^{s-1}$, the Motzkin numbers satisfy
$M_{mp^k+d}\equiv T_hM_d$ for $0\le d\le q-3$, with explicit corrections at
$d=q-2$ and $d=q-1$; the boundary case gives
$M_{np^r-2}\equiv\tfrac{\varepsilon_p^r-1}{2}T_{n-1}-U_{n-1}\pmod p$.  Adds the
exact first failures beyond the two initial-block ranges, arbitrary-multiplier
and all-prime-power boundary extensions, a prime-modulus digit algorithm, and
examples showing why unweighted repetition modulo a prime cube is generally
false.  Discusses Burns, Kohen and Pan--Sun rather than claiming novelty.}

\subsection{Congruences for iterated and compositional series}

\entry{Arithmetic of a Self-Iterating Generating Function}
{congruences-and-valuations/iterated-series/a396807-self-iterating-series}
{article.pdf}{A396807\_congruences\_and\_prime\_power\_lifts.zip}
{Both conjectures in OEIS A396807 for $A(x)=x+A^{[5]}(x)A^{[6]}(x)$ (functional
iteration): $A(x)\equiv x/(1-x)$ and $A^{[k]}(x)\equiv x/(1-kx)$ modulo $10$ for
every integer $k$, negative iterates included.  Adds a sharp modulus theorem
for $F=x+cF^{[r]}F^{[s]}$, rational lifts at every selected prime power, and
minimal period $46860$ for A396807 modulo $100$ starting at index~1.}

\entry{Prime-power congruences for compositional tree series}
{congruences-and-valuations/iterated-series/compositional-tree-series-congruences}
{article.pdf}{compositional\_tree\_congruences.zip}
{For $A_l(x)=x\exp(A_l^{\circ l})$, shows the OEIS A396805 claim
$a_5(n)\equiv n\ (\mathrm{mod}\ 5)$ is \emph{false}: first failure $n=23$, with
exception set $\{20j+3,\,20j+4: j\ge1\}$ and residue~2 throughout.  Proves
eleven other conjectural assertions in A396803--A396806 via a general reduction
of coefficients mod $p^\alpha$ to typed rooted trees on fewer than $p^\alpha$
labels, yielding eventual periods, rational residue generating functions, and
an exact classification of $a_l(n)\equiv n\pmod p$.}

\subsection{Valuations, divisibility and periodicity}

\entry{Binary Carries and the Divisibility of Two-Stack-Sortable Permutation Numbers}
{congruences-and-valuations/a000139-binary-carries-parity}
{research.pdf}{a000139\_carries\_research.zip}
{Proves Peter Bala's parity conjecture for OEIS A000139,
$a(n)=2(3n)!/((n+1)!(2n+1)!)$: $a(n)$ is odd exactly when $n$ is odd and its
binary expansion contains no \texttt{11}.  The full valuation is
$v_2(a(n))=s(n)+s(n+1)-s(3n)$ in binary digit sums.  Two independent parity
proofs, plus exact valuation distributions, recurrences, asymptotics, extremal
formulas and limit laws.}

\entry{Polynomial divisibility of height-one factorial ratios}
{congruences-and-valuations/factorial-ratio-polynomial-divisibility}
{article.pdf}{bala\_factorial\_divisibility\_research.zip}
{A general rational-root criterion for uniform polynomial divisibility of
balanced integral factorial ratios, applied to Peter Bala's uniform-product
conjectures for OEIS A211417 and A295431: explicit multipliers for every product
length, with the least multipliers certified in eight particular examples.
Includes an atlas of the 52 sporadic height-one ratios with recurrences,
generating functions and asymptotic expansions.}

\entry{When the Secant Numbers Become Periodic}
{congruences-and-valuations/secant-number-periodicity}
{article.pdf}{secant\_periodicity\_research.zip}
{Refutes the OEIS A000364 conjecture that the secant numbers from $n=1$ are
purely periodic modulo every $m$: at $m=27$, $a(1)\equiv1$ but
$a(19)\equiv10$.  The exact replacement is that pure periodicity from index one
holds if and only if no odd prime cube divides $m$; the least eventual period
always divides $\varphi(m)$.  Determines the least eventual period and exact
onset for every modulus, with densities and a certified mean onset.}

\section{Tetration and digit stabilization}

\entry{A Counterexample to a Decimal-Length Bound for Tetration Stabilization}
{tetration-and-digit-stabilization/congruence-speed-decimal-length-bound}
{article.pdf}{tetration\_counterexample.zip}
{An explicit 2544-digit base $A=1+qu$ ($q=5^{2547}$) for which
$\nu_{10}(T_{n+1}-T_n)=\min(1+2548n,\,2547(n+1))$ at every height, so the
congruence speed at height $\mathrm{len}(A)+2$ is $2548$ while the permanent
value is $2547$.  This contradicts the decimal-length form of Conjecture~1 of
Marco Rip\`a, \emph{The congruence speed formula} (NNTDM 27(4), 2021).  The
report audits the later Rip\`a--Onnis confirmation: their weaker valuation
bound is compatible with the counterexample, the decimal-length bound is not.}

\entry{Tetration, perfect powers, and frozen decimal digits}
{tetration-and-digit-stabilization/perfect-power-minima-for-frozen-digits}
{article.pdf}{tetration\_perfect\_power\_minima.zip}
{Determines $\rho(n)$, the smallest $n$th perfect power $x^n>1$ whose eventual
number of newly frozen terminal decimal digits is exactly $n$:
$\rho(1)=2$, $\rho(2)=49$, and $\rho(n)=B(s)^n$ for $n\ge3$ with
$s=n-v_2(n)$ and $B(s)\in\{2^s\pm1,\,3\cdot2^s\pm1\}$ according to
$s\bmod4$.  All other roots are ruled out by a deterministic quadratic size
barrier plus small exact certificates.}

\entry{Phase-shift classification for integer tetration}
{tetration-and-digit-stabilization/phase-shift-classification}
{tetration\_phase\_classification.pdf}{tetration\_phase\_classification.zip}
{Proves the 21-word phase-shift conjecture in OEIS A376842, the 14-word
even-base restriction from the appendix of Rip\`a's 2025 paper, and the
28-word prescribed-anchor conjecture in OEIS A376446.  The mechanism compares
the 2-adic and 5-adic valuations of $D_b=T_{b+1}-T_b$ at the first
permanently constant-speed height: strict 2-adic dominance gives phase 5,
strict 5-adic dominance a multiplicative orbit in $\mathbb{F}_5^\times$, and
balanced valuations exactly four further anchored words.}

\entry{A Sharp Convergence Theorem for the $q$-Newton Series of Subcritical Tetration}
{tetration-and-digit-stabilization/q-newton-series-convergence}
{article/tetration\_qnewton.pdf}{tetration\_qnewton\_research.zip}
{Proves convergence of the explicit $q$-binomial interpolation series of
MathOverflow question 259278 for every real base $1<a<e^{1/e}$, with
$\log a=qe^{-q}$: the grouped series converges absolutely and locally uniformly
for $\mathrm{Re}\,z>-2$, converges conditionally on $\mathrm{Re}\,z=-2$ except
at $z=-2+2\pi ik/(-\log q)$, and diverges for $\mathrm{Re}\,z<-2$.  Adds strict
complete monotonicity and uniqueness under an eventual log-convexity
condition.}

\entry{Stabilization of integer power towers}
{tetration-and-digit-stabilization/power-tower-stabilization}
{article.pdf}{tetration\_stabilization\_research.zip + tetration\_exact\_digits.zip}
{Refutes the claim that every prime base ending in 9 reaches its permanent
decimal congruence speed by height~2 (Rip\`a, Conjecture~2; arXiv:2208.02622v1,
Conjecture~4.1): $2749$ is the smallest prime counterexample, with
$H(2749)=3$.  Proves exact formulas for $C_a(h)$, $V(a)$ and $H(a)$, an exact
natural density for onsets, and infinitely many primes at each prescribed onset
via CRT and Dirichlet.  Also resolves the three conjectural comments in OEIS
A324017 --- two proved, the third disproved, with
$3\uparrow\uparrow2\equiv27$ and $3\uparrow\uparrow3\equiv59\pmod{64}$ --- via
$T_{h+r}-T_h\equiv-2q^h\pmod{q^{h+1}}$ for even $q\ge4$, so exactly one radix-$q$
digit freezes per height.  Adds a factorization-free digit-lifting algorithm,
modular saturation across Knuth arrow ranks, and a sign-sensitive local $p$-adic
Lambert $W$ formula.  \emph{Merged; the shared one-digit-per-height theorem was
proved identically in both originals and is proved once here.}}

\entry{Delayed digit stabilization: a counterexample and sharp bounds}
{tetration-and-digit-stabilization/iterated-power-digit-stabilization}
{delayed\_digit\_stabilization.pdf}{digit\_stabilization\_research.zip}
{For $S_a(b)=v_{10}(a^{10^{b+1}}-a^{10^b})$ and $D_a(b)=S_a(b)-S_a(b-1)$,
answers negatively the question whether $D_a(b)=1$ for all $b\ge3$: the
smallest counterexample is $a=3\cdot2^{99}+1$, with $S_a(2)=100$,
$S_a(3)=102$, $D_a(3)=2$.  Proves minimality, an exact onset formula, the
smallest counterexample at every stage, the natural density of exceptions, and
arbitrary-radix extensions.  The article is explicit that $a^{10^b}$ is not
genuine tetration and derives the tower formulas separately.}

\section{Log-concavity, unimodality and positivity}

\entry{Single curves do not force log-concavity}
{log-concavity-and-unimodality/cluster-variable-log-concavity-counterexample}
{article.pdf}{annular\_single\_lamination\_counterexample.zip}
{An unpunctured annulus with one marked point on each boundary and a single
unit-weight simple lamination curve gives, after mutations $2,1,2$ with the two
mutable variables specialised to~1, the polynomial
$1+3q^2+2q^3+4q^4+2q^5+q^6$; the consecutive coefficients $3,2,4$ form a strict
valley, refuting Conjecture~1 of arXiv:2508.04396v3 and conflicting with its
general unpunctured-surface unimodality statement.  The geometric realisation
is proved, not merely the matrix calculation, and a fixed seed gives infinitely
many such examples.}

\entry{Cycles refute infinite log-concavity for chromatic coefficients}
{log-concavity-and-unimodality/infinite-log-concavity-of-cycles}
{article.pdf}{chromatic\_cycle\_logconcavity.zip}
{For the cycle $C_n$ with absolute chromatic coefficients under the iterated
operator $L(a)_k=a_k^2-a_{k-1}a_{k+1}$: infinitely log-concave for
$3\le n\le11$; first failure at depth 5 for $n=12$, depth 4 for
$13\le n\le16$, and depth 3 for $n\ge17$.  So $C_{12}$ is the smallest cycle
counterexample (no minimum over all graphs is claimed), with an explicit
third-iterate witness polynomial valid for all $n\ge3$.}

\entry{Oscillating Log-Concavity in Geometric Generating Functions}
{log-concavity-and-unimodality/oscillating-log-concavity-geometric-transform}
{article.pdf}{oscillating\_log\_concavity\_research.zip}
{Challenges 2 and 3 in Section~3 of Heim and Neuhauser, arXiv:2302.13327v1, for
the \emph{geometric} transform $1/(1-uG(z))$.  For divisor-sum weights, every
real exponent and every $u>0$ give infinitely many strict log-concavity and
log-convexity indices with eventually bounded gaps; likewise for integer power
weights with exponent ${\ge}\,2$, the exponents $0$ and $1$ being explicitly
exceptional.  For square weights at $u=1$, every four consecutive indices carry
both signs (sharp), each with density $1/2$ in every arithmetic progression.
A subdominant-pole theorem is the common mechanism.}

\entry{A counterexample and a sharp spacing-three criterion for products of $q$-integers}
{log-concavity-and-unimodality/q-integer-product-unimodality}
{report.pdf}{q\_unimodality\_attack.zip}
{The necessity half of Conjecture~5.4 of Connelly, Ito, Martinez, Shevchenko and
Yang (arXiv:2605.12822v1) is false: $(1+q)^6(1+q^3)$ is symmetric and weakly
unimodal although no $a_i$ is divisible by 3 and $b=2$ exceeds the asserted
bound, and the example lies inside the source's claimed range.  The report also
proves the proposed sufficient condition for every $r\ge2$, a necessary degree
bound in the absence of an $r$-divisible ordinary factor, and a complete
classification of spacings two and three.}

\entry{Sharp size thresholds for failures of log-concavity in independence systems}
{log-concavity-and-unimodality/independence-system-thresholds}
{article.pdf}{sharp\_log\_concavity\_thresholds.zip}
{For each rank $r\ge4$, determines the smallest ambient ground-set size
admitting failure of ordinary, ordered and ultra log-concavity, and shows the
three thresholds coincide across all independence systems, those satisfying
Xie--Xu's augmentation bound, and those with additive augmentation parameter~3
(values $7/6/5$, $7/7/6$, $8/7/7$, and $r+1$ throughout for $r\ge7$).  The
initial target, Xie--Xu Conjecture~4.8 (arXiv:2408.09152v1), had already been
disproved by Alex Chengyu Li in September 2026; the report states this and
refines the target to the rankwise minimum-size question.}

\entry{A sharp order-five threshold for higher-order Stirling subset log-concavity}
{log-concavity-and-unimodality/higher-order-stirling-rows}
{article.pdf}{stirling\_log\_concavity.zip}
{For $S^{(r)}_{n,k}$, the number of partitions of $n+(r-1)k$ labels into $k$
blocks of size ${\ge}\,r$, proves that all rows are log-concave if and only if
$1\le r\le5$, with strict inequalities in the interior --- the log-concavity
assertion of Deb--Sokal Conjecture~1.4(c) (arXiv:2507.18959v1).  Also proves the
cycle cases $r=3,4$; the fifth-order \emph{cycle} conjecture is explicitly left
unresolved, and the included operator counterexample is flagged as not being a
counterexample to it.}

\entry{A sum-of-squares proof of Ulas's real-zero conjecture}
{log-concavity-and-unimodality/stern-polynomial-positivity}
{stern\_conjecture\_note.pdf}{stern\_conjecture\_research.zip}
{Conjecture~4.4 of Maciej Ulas, arXiv:1909.10844v1, for the Stern polynomials
indexed by $a_n=(4^n-1)/3$ and $h_n=2(4^n-1)(2\cdot4^n+1)/3+1$: for every
$n\ge0$ and every real $t$,
$W_n(t)=(A_{n+1}(t)-tA_n(t)/2)^2+3t^2A_n(t)^2/4>0$.  The adjacent-polynomial
quadratic identity behind it was already in Ulas's paper and is re-derived, not
claimed as new.}

\entry{Two-coin counterexamples to generalized Shepp--Olkin thresholds}
{log-concavity-and-unimodality/bernoulli-entropy-quasiconcavity}
{article.pdf}{bernoulli\_entropy\_research.zip}
{For every finite real $q>1$, both the R\'enyi-$q$ and the Tsallis-$q$ entropy of
a sum of independent Bernoulli variables fail even \emph{quasi}concavity as a
function of the parameter vector, already on the open square $(0,1)^2$ with two
coins.  This disproves the proposed thresholds --- R\'enyi~2 and Tsallis
$3.65986\ldots$ --- adjoined to Conjecture~4.2 of arXiv:1503.01570 and restated
in arXiv:1510.05390, Section~8.}

\entry{Nonreal roots survive in polynomial iterative equations}
{log-concavity-and-unimodality/nonreal-roots-in-iterative-equations}
{article.pdf}{nonreal\_roots\_research.zip}
{Here $P[f]=\sum_ja_jf^{\circ j}$ means a combination of \emph{iterates}, not a
polynomial in pointwise values.  Gives a negative answer to Problem~6.1 of
Draga--Morawiec (arXiv:1503.00570), which proposed deleting nonreal
maximal-modulus characteristic roots: an explicit odd, strictly increasing,
globally bi-Lipschitz homeomorphism realises a cubic whose nonreal root cannot
be discarded.  Adds a classification of the circular root spectra that occur.}

\entry{Canonical binomial decomposition does not preserve real roots}
{log-concavity-and-unimodality/binomial-decomposition-real-roots}
{article.pdf}{mu\_welker\_recursive\_decomposition.zip + recursive\_binomial\_counterexample.zip}
{A negative answer to Question~3.10 of Mu and Welker (arXiv:2503.24076v1) for
their canonical decomposition $F(t)=G(t)+tH(t)$ built from the integer binomial
expansion of every coefficient: real-rootedness of $F$ does not pass to the
parts.  Two cubic witnesses are carried, one with a repeated root and one with
simple roots, each with its own digit table, discriminant pair and exact
nonreal zeros; the report also gives the minimal witness within cubics, two
all-degree thresholds neither of which implies the other, and a two-sided
limit.  \emph{Merged; the two originals were two write-ups of one proof
program, each already containing both witnesses with the headline and the
co-star swapped.}}

\entry{A first-triple criterion for log-concavity of MDS weight distributions}
{log-concavity-and-unimodality/mds-weight-log-concavity}
{article.pdf}{mds\_log\_concavity\_research.zip}
{For $d=n-k+1\ge2$, $k\ge3$ and $q>d$, the compressed formal MDS weight
distribution $(1,A_d,\ldots,A_n)$ is log-concave if and only if the single
inequality $A_{d+1}^2\ge A_dA_{d+2}$ holds --- the whole sequence is controlled
by its first triple.}

\entry{A uniform Hankel obstruction for higher-order Stirling subset polynomials}
{log-concavity-and-unimodality/stirling-hankel-obstruction}
{article.pdf}{stirling\_hankel\_obstruction.zip}
{Proves the higher-order negative assertion of Deb--Sokal Conjecture~1.4(d),
with a stronger statement at every positive shift: for $r\ge3$ and $a\ge1$ the
$3\times3$ Hankel determinant $\det(s_{r,a+i+j}(x))$ of the row-generating
polynomials fails the required positivity.  Clause~(d) of that conjecture;
clause~(c) is the subject of the higher-order Stirling rows entry.}

\section{Quaternionic analysis}

One report, merged from ten manuscripts.  The subject is new to this
collection, and the delivery was unusual: one expository survey and nine
independent attempts at the same follow-up problem, all written on the same
day.  The nine agree, which is the substance of the entry.

\entry{Slice regularity and the global inverse Fueter problem}
{quaternionic-analysis/slice-regularity-and-fueter-inversion}
{article.pdf}{quaternionic\_analysis.zip and nine companion archives}
{One article assembled from ten manuscripts.  Part~I is the delivered
expository survey: slice regularity and Cauchy--Fueter regularity side by side,
with Cauchy formulas, Taylor and Laurent expansions, the regular product and
the geometry of zeros, the argument principle, Rouch\'e, open mappings,
Schwarz--Pick and an intrinsic Riemann mapping theorem.  Parts~II--XI answer
the follow-up question the other nine manuscripts all attacked independently:
for an axially monogenic $g=P+IQ$, when is there a single-valued slice-regular
$f$ with $\Delta_4 f=g$?  The answer is a two-period obstruction.  The closed
one-forms $\omega_g=\tfrac12(-P\,dx+Q\,dy)$ and
$\theta_g=\tfrac12((xP+yQ)dx+(yP-xQ)dy)$ have vanishing periods exactly when a
global primitive exists, and then $f(x+Iy)=(x+Iy)C+V$ with $dC=\omega_g$,
$dV=\theta_g$, unique modulo $q\mapsto qa+b$; the monodromy is affine, and the
cokernel of the Fueter map has real dimension $8m$.  The point the nine exist
to make is that $\theta_g$ is computed from the target, with no first primitive
chosen.  \textbf{Nine independent derivations reaching the same two forms, the
same $8m$ and the same logarithmic energy constant is the evidence the entry
rests on.}  Where they differed it was almost always notation --- thirteen
conflicts are recorded and resolved in \nolinkurl{MERGE_NOTES.md} --- but two
were mathematical and would have produced a false statement if the merge had
been mechanical: the size law at a critical singularity has genuinely different
sharp constants in two norms that four manuscripts both called $M_g$, and the
conserved flux recovering the slope residue needs the weight
$3(\operatorname{Re}q-u)+\operatorname{Im}q$ rather than $q-u$, the coefficient $3$ being exactly what
makes it right-Fueter-regular.  All ten source manuscripts and all ten
verification programs are kept; the deduplication is in the article.}

\section{Graph theory}

\entry{Mutual-visibility sets in Andr\'asfai graphs}
{graph-theory/andrasfai-mutual-visibility}
{article.pdf}{andrasfai\_visibility\_research.zip}
{Proves the generating function attributed to Andrew Howroyd in OEIS A391632,
$\sum a(n)x^n = x(4-19x+33x^2-40x^3+8x^4)/(1-10x+29x^2-32x^3+8x^4)$, by an
exact weight-preserving bijection between mutual-visibility sets of the
$n$-Andr\'asfai graph and anchored closed walks of length $3n-1$ in a
four-state labelled graph.  Adds the full bivariate generating function, an
exact binomial formula for every cardinality, maximum-set classification,
eventual polynomiality at each fixed deficit, exponential asymptotics, a
central limit theorem, and linear-time recognition and exact uniform sampling.}

\entry{A bicyclic 30-vertex graph with domination root $-4$}
{graph-theory/domination-root-minus-four-order-30}
{domination\_root\_30.pdf}{domination\_root\_30\_research.zip}
{An explicit connected planar graph on 30 vertices and 31 edges, treewidth 2 and
domination number 10, with $D(G,x)=x^{10}(x+4)Q_{19}(x)$ --- a negative answer
to question~2 of Section~4 in Alikhani and Griswold, arXiv:2608.00109v1, which
asks whether 33 is the minimum order.  Their 33-vertex example had already
disproved the original integer-root conjecture; no claim of global minimality
at 30 is made.  Three independent exact polynomial computations agree.}

\entry{A marked-cycle proof of the trapezohedral domination conjecture}
{graph-theory/trapezohedral-minimal-dominating-sets}
{article.pdf}{trapezohedral\_domination\_research.zip, A381190\_proof\_package.zip}
{Proves the generating function labelled conjectural in OEIS A381190: for the
$n$-trapezohedral graph, the number of dominating sets that are
inclusion-minimal among all dominating sets and induce a connected subgraph is
$a_n=2n(2p_{n-2}+p_{n-3})$ with $p_m=p_{m-2}+p_{m-3}$ (Padovan).  The proof
classifies such sets by a cyclic arrangement of gaps of length 2 and 3 with a
restricted marked edge.  Adds the size distribution, strict log-concavity, an
unbiased sampler, a Binet formula and a central limit theorem.}

\section{Enumerative combinatorics}

\entry{A Twelve-Value Proof of the A181280 Formula}
{enumerative-combinatorics/a181280-binary-matrix-formula}
{article.pdf}{A181280\_research\_package.zip}
{The binary-matrix formula recorded as Conjecture~20 by Kauers and Koutschan
(2023) and still marked conjectural in OEIS A181280: an explicit spectral
recurrence bound and a finite certificate resting on twelve decisive counts,
together with further consequences.  \texttt{STATUS.md} records that general
C-finiteness was already proved in Kauers's December 2025 lecture slides.}

\entry{Esthetic Numbers, Folded Paths, and Algebraic Diagonals}
{enumerative-combinatorics/esthetic-numbers-and-folded-paths}
{article.pdf}{esthetic\_walks\_research.zip}
{The diagonal conjecture $T(q,q-2)=A182555(q-2)$ in OEIS A377000, where
$T(q,k)$ counts base-$q$ positive integers of length $k$ whose adjacent digits
differ by one.  Proves it exactly, proves the entry's other two conjectures as
retrieved, and derives formulas for every fixed-offset diagonal $q=k+d$.}

\entry{Diagonal Hardinian Arrays: A Fixed-Parameter Asymptotic Formula}
{enumerative-combinatorics/hardinian-array-diagonal-asymptotics}
{article.pdf}{hardinian\_arrays\_research.zip}
{For every fixed $r\ge0$, $H_r(n,n)\sim C_r4^{rn}n^{-r(r-1)/2}$ with an explicit
$C_r$ in Gamma values.  The formula confirms the constants tabulated by
Dougherty-Bliss and Kauers (2024) but refutes their proposed restriction to
powers of 2, 3 and $\pi$: $C_7=2^{18}5^2/(3^{45}\pi^3)$.  Adds an exact finite
determinant at fixed array size, an exact enumeration algorithm, and a limiting
negative-binomial law for total deficits.}

\entry{An explicit formula for lucky parking spots}
{enumerative-combinatorics/lucky-parking-spots}
{article.pdf}{lucky\_parking\_spots\_research.zip}
{Proves Conjecture~13 of Butler, Hadaway, Lenius, Martens and Moats (JIS 29
(2026), Article 26.1.1): a closed form for $C(n,j)$, the number of parking
functions in which \emph{spot} $j$ is lucky.  Adds an exact reflection law, all
fixed right-edge columns, a closed form for OEIS A374533, Poisson edge limits,
uniform bulk convergence, polynomial reciprocity and asymptotic expansions.}

\entry{A hexagonal-lattice proof of the A195806 conjecture}
{enumerative-combinatorics/a195806-hexagonal-lattice}
{article.pdf}{A195806\_hexagonal\_proof.zip}
{An elementary derivation of the six residue-class formulas and the
constant-coefficient recurrence recorded as conjectural in OEIS A195806,
starting from the original triangular-array constraints rather than from
interpolation.  The steps are an integral six-parameter bijection, four
independent translation intervals, a dihedral chamber reduction, and an exact
lattice count.}

\entry{Leaves of large type in the numerical-semigroup tree}
{enumerative-combinatorics/numerical-semigroup-leaf-types}
{article.pdf}{leaf\_type\_counterexample.zip}
{Disproves Conjecture~4.4 of Chappelon, Ram\'irez Alfons\'in and Stamate,
\emph{Numerical semigroup tree: type-representation} (arXiv:2507.15006v2), with
an explicit genus-34 counterexample, and establishes a sharp asymptotic order
for the type of leaves in that tree.}

\entry{Strict growth and the Catalan limit for adjacency-bounded 132-avoiding permutations}
{enumerative-combinatorics/adjacency-bounded-132-avoiders}
{article.pdf}{strict\_growth\_132.zip}
{For $132$-avoiding permutations with every adjacent difference at most $m$ and
growth constant $\alpha_m$: proves $\alpha_{m+1}>\alpha_m$ for all $m\ge1$ ---
the remaining strict-inequality part of Nadler's Conjecture~2 --- and the
refinement $4-\alpha_m=(2\log m+4\log\log m+O(1))/m$, with
$2\log\pi-4\log2\le\liminf E_m\le\limsup E_m\le2\log\pi$.  The all-$m$
component ordering and convergence of $E_m$ are explicitly not claimed.}

\entry{Throwback Dynamics with Repeated Entries}
{enumerative-combinatorics/throwback-queue-dynamics}
{throwback\_research.pdf}{throwback\_research.zip}
{For the queue process that removes the head and reinserts it at the zero-based
position equal to its weight: every token recurs if and only if at most $m$
input tokens have weight at most $m$, for every $m\ge1$, with an explicit
description of the recurrent tokens when that fails.  Adds exact cycle periods
and frequencies, a coprimality criterion for perfect mixing, Catalan and
prime-parking counts for recurrent cores, and the largest labelled period of a
core of size $r$.}

\entry{Reciprocity and Matrix Duality for Preorder Polytopes}
{enumerative-combinatorics/preorder-polytope-reciprocity}
{article.pdf}{preorder\_reciprocity.zip}
{Proves Conjectures 7.2, 4.2 and 4.4 of Athanasiadis and Chapoton,
\emph{Polytopes and posets associated to preorders} (arXiv:2605.26916v1), and
the ordinary $q=1$ case of their Conjecture~4.9 (not its full $q$-refinement).
Two mechanisms: Ehrhart--Macdonald reciprocity after subtracting the all-ones
vector, extending the double identity to independently weighted capacities; and
duplication of both colour classes of the bipartite graph followed by
Postnikov's lattice-point formula, giving the conjectured matrix transpose.}

\entry{Support reciprocity and the $q$-zeta polynomial of preorder lattice-point posets}
{enumerative-combinatorics/preorder-q-zeta-reciprocity}
{article.pdf}{preorder\_qzeta\_reciprocity.zip + preorder\_q\_zeta\_reciprocity.zip + preorder\_q\_reciprocity\_research.zip}
{Proves the full $q$-refinement of Athanasiadis--Chapoton Conjecture~4.9,
equation~(21): for a preorder $\tau$ on $[n]$, with $P_\tau$ the nonnegative
integer vectors satisfying $a(I)\le|I|$ on every order ideal and $B_\tau$ its
maximal elements, $Z_q(P_\tau,[-1]_q)=(-1)^n\sum_{b\in B_\tau}q^{-|\mathrm{supp}(b)|}$,
together with the stronger identity at each exact support.  The $q=1$ case was
already covered by the preorder-polytope reciprocity report and is cited, not
reclaimed.  \emph{Merged from three reports that reached the same two theorems
by three architecturally independent routes --- an incidence inversion with
Boolean localization, the same with a Hall-based privatization and a
self-contained Euler identity, and a capacity-enumerator reciprocity with a
transfer involution --- and all three survive.}}

\entry{An explicit shelling for preorder lattice-point posets}
{enumerative-combinatorics/preorder-shellability}
{article.pdf}{preorder\_shellability\_research.zip}
{An affirmative solution to Question~4.6 of Athanasiadis and Chapoton
(arXiv:2605.26916v1, p.~13): an explicit pure shelling of the order complex of
$P_\tau$ with its minimum removed, and a coning that shells the order complex
of $P_\tau$ itself.  Both complexes are therefore Cohen--Macaulay over the
integers and over every field.  A different question of that paper from the
three other preorder entries.}

\entry{Multiple-chain preorder polytopes}
{enumerative-combinatorics/multiple-chain-exponential-formula}
{article.pdf}{multiple\_chain\_exponential\_formula\_research.zip + multiple\_chains\_research.zip}
{Proves the exponential generating-function identity proposed in Section~9.5 of
Athanasiadis and Chapoton (arXiv:2605.26916v1) for lattice points of
multiple-chain preorder polytopes counted by number of nonzero coordinates, in
a rectangular generalization with independent block width and height.  Adds a
duality corollary, an Ehrhart evaluation, a degree-$2^d$ root theorem with a
quartic elimination certificate, two treatments of the $\gamma$-vector, a
Stieltjes moment representation, the boundary constant in two forms, and a
central limit theorem with exact variance.  \emph{Merged; the identity is
proved twice, once through a subset encoding with a bounded symbol and once
through slack meanders in a Laurent ring, and the eleven deliberate
redundancies are listed in the text.}}

\entry{Rationality and growth of stabilized permutation-weight series}
{enumerative-combinatorics/stabilized-permutation-weights}
{article.pdf}{stabilized\_permutation\_weights.zip}
{Studies the full stabilized generating functions of the permutation-weight
$q$-Eulerian polynomials, beyond their initial connection with OEIS A256193:
an all-$d$ finite-spectrum theorem, rationality, a universal integer
denominator, the exact radius of convergence and residue-class asymptotics.
The two-descent sequence gets an explicit formula in powers of 3, powers of 2
and Fibonacci numbers.}

\entry{A Reflexive Root-Polytope Model for Preorder $h$-Polynomials}
{enumerative-combinatorics/preorder-root-polytopes}
{article.pdf}{preorder\_root\_polytopes.zip}
{Proves Conjecture~5.1(c) of Athanasiadis and Chapoton (arXiv:2605.26916v1): for
a preorder $\tau$ on $[n]$ the directed root polytope
$C_\tau=\mathrm{conv}(\{\pm e_i\}\cup\{e_i-e_j: j\le_\tau i\})$ is reflexive with
$h^*(C_\tau,t)=h_\tau(t)$, and any pulling triangulation of its boundary is
combinatorially the boundary of a simplicial polytope with that $h$-polynomial.
Palindromicity and unimodality, Conjecture~5.1(a,b), follow, as do the full
$g$-theorem conditions.  The higher-dimensional augmented bipartite identity and
preorder duality are credited to Dai, Hou, Liu, Thawinrak and Wang
(arXiv:2608.16037v2); flag realizability, $\gamma$-positivity and real-rootedness
are not proved.  A different conjecture of the same paper from the preorder-polytope
reciprocity report.}

\entry{Largest Simplex Faces under Ordinal Sums}
{enumerative-combinatorics/simplex-faces-of-order-and-chain-polytopes}
{Largest\_Simplex\_Faces.pdf}{Series\_Parallel\_Simplex\_Faces\_Research.zip}
{Answers the concluding open question of Freij-Hollanti and Lundstr\"om,
\emph{Simplex inequalities of order and chain polytopes of recursively defined
posets} (Order 43 (2026), article~31).  For $s_O(P)$ and $s_C(P)$, the largest
dimensions of simplex \emph{faces} of the order and chain polytopes: both are
computable in $O(n)$ on any series-parallel decomposition, the order case
carrying eight normalized endpoint states.  For an ordinal sum of antichains,
$s_O=r+\nu$ with $\nu$ a maximum $b$-matching of the layer path under capacities
$\min(a_i-1,2)$, and $s_C=r+\max(0,k-1)$.  The extremal gap is sharp:
$0\le s_C-s_O\le\lfloor(n-2)/3\rfloor$, attained for every $n\ge2$.}

\entry{Skew partitions, continued fractions and connected components}
{enumerative-combinatorics/skew-partition-continued-fractions}
{article.pdf}{A225114\_proof\_bundle.zip}
{Proves the repeated-power continued-fraction formula recorded as conjectural in
OEIS A225114, which counts skew partition diagrams with no empty row or column.
Unique concatenation gives $A(q)=1/(1-P(q))$ for the edge-connected diagrams,
and a boundary-path encoding followed by a weighted two-coloured Motzkin-path
decomposition gives
$P(q)=q/(1-2q-q^3/(1-2q^2-q^5/(1-2q^3-\cdots)))$, whence $A(q)=1/(2-C(q))$ with
$C=1+P$.  Adds a four-variable refinement, a formula by height, width and
component count, a sharp truncation bound, exponential asymptotics and a central
limit theorem for the number of components.  \textbf{Explicitly not a novelty
claim}: a prior-art check found Pasechnik's MathOverflow sketch of ten September
2024, and the entry credits Kurkov; the report says so first and credits both.}

\entry{From mesh avoidance to Catalan inflation}
{enumerative-combinatorics/mesh-avoidance-catalan-inflation}
{article.pdf}{A289587\_proof.zip, OEIS\_A289587\_proof.zip}
{Proves Scheuerle's generating-function conjecture in OEIS A289587 for
permutations avoiding $321$ and the mesh pattern $(12,174)$, and the equal count
for $(12,234)$ by inversion: with $D=x^4-2x^3-5x^2-2x+1$,
$P=x^4+4x^3+11x^2+10x+3$ and $Q=x^2+5x+3$, the generating function is
$(P-Q\sqrt D)/(8x(1+x)^2)$.  The headline structural result is the exact
\emph{occurrence count} $\mathrm{occ}=\binom{f}{2}+\sum_U\binom{|U|}{2}$, of
which avoidance is the case $\mathrm{occ}=0$.  Merged from two reports that
reached the identical generating function by the identical run-contraction
substitution $u\mapsto u/(1+xu)$; the merge keeps \emph{both} refinements,
because they refine the same sequence by different statistics --- excedances and
records --- and are genuinely different theorems, agreeing only in total.  The
two local predicates they are built on disagree on $31$ of the $321$-avoiders of
length at most six, the smallest being $12$, and the merged article proves where
they coincide rather than assuming it.}

\entry{Prescribed coefficient-lattice minima}
{enumerative-combinatorics/coefficient-lattice-minima}
{article.pdf}{coefficient\_lattice\_minima\_research.zip}
{Settles the unnumbered realization conjecture after Lemma~1 of O'Bryant,
\emph{On Nathanson's Triangular Number Phenomenon} (arXiv:2506.20836): every
nondecreasing list of even integers is the list of successive $L^1$-minima of
the coefficient lattice of some finite integer set.  The witness is the
mixed-radix set $A=\{0,1,h_1,h_1h_2,\ldots\}$, whose minima are exactly
$(2h_1,\ldots,2h_r)$ with $r+2$ distinct elements, and the radices may be given
in any order.}

\entry{Coefficientwise magic positivity through dimension ten}
{enumerative-combinatorics/magic-positivity-parking-polytopes}
{article.pdf}{magic\_positivity\_research.zip}
{A polytope is \emph{magic positive} when its Ehrhart polynomial has nonnegative
coefficients in the basis $t^j(t+1)^{n-j}$.  Proves, in exact arithmetic, that
the generalized parking-function polytope $X_n(b)$ is magic positive for every
dimension $3\le n\le10$ and \emph{every} positive integer parameter vector,
uniformly in $b$ rather than on a sampled box --- a partial solution of
Conjecture~8.1 of Hill, Luo, Trinh and Vindas-Mel\'endez (arXiv:2607.15503v1).
Dimension three is proved by printed algebraic formulas; the first-coefficient
harmonic-number formula and the endpoint identities hold in every dimension.
The conjecture in unrestricted dimension remains open.}

\entry{Monotone blocks and reverse replies}
{enumerative-combinatorics/reverse-reply-avoidance-game}
{article.pdf}{reverse\_reply\_solution.zip}
{Proves Conjecture~2.9 of Ulfarsson (arXiv:2603.16004): for every $k\ge3$ and
every $n\ge R(k)=k^4-2k^3-3k^2+5k+5$, the second player wins the normal-play
permutation-avoidance game on $S_n$ with pattern length $k$ by always replying
with the reversal, so the Sprague--Grundy value is~0.  The proof is elementary,
with no computer search: a minimal-support lemma isolates one pattern inside a
family built from singletons plus one inflatable monotone interval, and
inflation restores the original length so the reply is legal.  $R(k)$ is
asserted sufficient, not minimal.}

\entry{Nonintegral tropical degree from two horizontal segments}
{enumerative-combinatorics/tropical-hilbert-nonintegral-degree}
{article.pdf}{tropical\_hilbert\_counterexample.zip}
{Computes the strict-witness tropical Hilbert function of Grigoriev
(arXiv:2404.06440v2) exactly for a two-segment set:
$\mathrm{TH}_{V_d}(k)=k+\min\{k,\lceil k/d\rceil\}+1$.  Taking $d=3$ gives
tropical degree $4/3$ --- not an integer --- and a quasipolynomial of least
period~3, so eventual polynomiality fails.  Every rational in $(1,2]$ is
realized, denominators are unbounded, and for $d=p/q>1$ in lowest terms the
least eventual period is exactly $p$.  Both bounds are proved by hand.}

\section{Generating functions and asymptotics}

\subsection{Coefficient counts of binary products}

Three reports on counting coefficients of $\prod_i(1+x^{G_i})$.  The first two
answer, independently, the universal-rationality part of Richard Stanley's
MathOverflow question 431075 (``A conjectured rational generating function'',
2022), which had no posted answer when checked on 19 September 2026.  A third
answer used the same mixed-base construction and the same enumerative argument
as the first, and has been merged into it.

\entry{A Mixed-Base Counterexample to Stanley's Rationality Question}
{generating-functions-and-asymptotics/binary-product-coefficient-counts/mixed-base-counterexample}
{article.pdf}{stanley\_mixed\_base\_counterexample.zip + Stanley\_Rationality\_Counterexample.zip}
{With $G_{2j+1}=2^j$ and $G_{2j+2}=3^j$ the exponent generating function is
rational, while $\nu(2n)=2^{\,n-\lfloor(n+1)\alpha\rfloor}$ for
$\alpha=\log_3 2$ counts the coefficients equal to one in
$\prod_{i\le N}(1+x^{G_i})$ and its generating function is not rational, not
D-finite, not P-recursive, and has a natural boundary.  Classifies rationality
in the family interleaving $a^j$ and $q^j$ with digit polynomial
$1+y+\cdots+y^{a-1}$ for all $2\le a<q$ --- Stanley's general $r=a-1$ --- as
holding exactly when $q^u=a^v$.  Adds a finite-quotient rigidity theorem, an
abstract natural-boundary theorem for $\sum q^{\lfloor n\theta\rfloor}t^n$ at
every irrational $\theta$, the exact limit set of $a_n/\lambda^n$, and a
streaming recurrence making the point that non-P-recursive does not mean hard to
compute.  \emph{Merged; the enumerative half was the same argument in both
originals, but the two routes to non-rationality differ and both are kept ---
an elementary finite-state reduction mod $p$, and the rigidity route that
reaches the stronger transcendence conclusions.}}

\entry{Rational exponents, nonrational coefficient counts}
{generating-functions-and-asymptotics/binary-product-coefficient-counts/recursive-exponent-construction}
{research.pdf}{stanley\_counterexample.zip}
{A different negative answer, from the recursion $a_0=2$, $a_1=1$,
$a_n=a_{n-1}-2a_{n-2}$, with $b_n=T_{n-1}+2^n-1+a_n$, $T_n=T_{n-1}+b_n$,
$G_{2n-1}=2^{n-1}$ and $G_{2n}=b_n$: the counting series has a natural boundary
on the unit circle, hence is neither rational, algebraic nor D-finite.}

\entry{Coefficient residues in generalized Fibonacci products}
{generating-functions-and-asymptotics/binary-product-coefficient-counts/kbonacci-coefficient-residues}
{article.pdf}{stanley\_kbonacci\_research.zip + quasifibonacci\_parity\_research.zip}
{Settles Conjectures 6.2 and 6.3 of Stanley, \emph{Theorems and Conjectures on
Some Rational Generating Functions} (arXiv:2101.02131v3, pp.~23--24).  For
$k$-bonacci seeds each exceeding the sum of their predecessors, the number of
odd coefficients of $\prod_{i\le n}(1+x^{w_i})$ has generating function
$(1+2z^k)/(1-2z+2z^k-2z^{k+1})$, independently of the seeds; the proof builds
periodic signs of block length $k+1$ and block product $-1$ and shows every
signed finite product is flat.  \emph{Merged; the Conjecture 6.3 halves proved
the same statement on the same domain, and the surviving report adds
Conjecture~6.2 and the asymptotic analysis.}}

\subsection{Single-sequence OEIS studies}

Nine reports on the growth or the closed form of a single OEIS sequence.  Three
of them count the terms in the derivatives of a power tower --- of $x^{x^x}$,
$x^{x^2}$ and $x^x$ --- and reach three different answers by three different
recurrences; two more share a Bell-scale saddle-point method on unrelated
sequences.  Shared technique is not shared theorem, and none of these nine
duplicates another.

\entry{Bell-scale growth of OEIS A088714}
{generating-functions-and-asymptotics/oeis-sequence-asymptotics/a088714-bell-scale-growth}
{article.pdf}{A088714\_asymptotic\_growth.zip}
{For the unique formal series with $A(x)=1+xA(x)^2A(xA(x))$, proves
$\log a_n=n\log n-n\log\log n-n+O(n\log\log n/\log n)$, so
$(e\log n/n)\,a_n^{1/n}\to1$ and the ordinary generating function has radius of
convergence zero while the exponential one is entire.  The lower bound comes
from a single explicit path through an exact triangular recurrence, the upper
from a Stirling-positivity comparison with derivatives of Touchard
polynomials.  The multiplicative equivalent is \emph{not} claimed, and the
finer Bell-number form is labelled a conjecture with a sufficient ratio
condition isolated.}

\entry{The exponential growth constant of OEIS A158415}
{generating-functions-and-asymptotics/oeis-sequence-asymptotics/a158415-growth-constant}
{article.pdf}{A158415\_asymptotic\_growth.zip}
{For the number of distinct values expressible with $n$ symbols from
$E::=1\mid\sqrt E\mid(E+E)$, proves a computable $\gamma>1$ exists with
$a(n)^{1/n}\to\gamma$ and $1.86204<\gamma<1.9208$, gives an effective
characterisation as a supremum of algebraic finite-data bounds, and proves the
boundary identity $\sum a(n)\gamma^{-n}=\gamma^2$, which rules out a pure
$C\gamma^n$ prefactor.  The rigorous lower bounds rest on a uniform injectivity
theorem: any finite alphabet of representable values can serve as nested-radical
digits, via an odd valuation at a prime above~2.  No closed form for $\gamma$ is
claimed, and terms 22--28 get certified intervals rather than exact values.}

\entry{An exact coefficient formula for Ramanujan's theta function}
{generating-functions-and-asymptotics/oeis-sequence-asymptotics/a260306-formula-proof}
{a260306\_proof.pdf}{A260306\_formula\_proof.zip}
{Proves both finite formulas recorded in OEIS A260306, for every coefficient
index, for the expansion $\theta(N)\sim\sum c_nN^{-n}$ defined by
$\sum_{k<N}N^k/k!+\theta(N)N^N/N!=e^N/2$.  The route is an exact Laplace-integral
difference, an analytic change of variable with $F(\psi(w))=w^2/2$, and Lagrange
inversion reducing each $c_n$ to one coefficient extraction, then a finite
change of basis into the entry's double and triple sums.  Includes the analytic
error estimate the identification needs, rather than only checking the two sums
agree, and identifies A065973 as the reduced denominators.}

\entry{The maximal positive coefficient of a Hermite polynomial}
{generating-functions-and-asymptotics/oeis-sequence-asymptotics/a277280-hermite-maximum}
{article.pdf}{A277280\_asymptotics.zip}
{An exact closed form for OEIS A277280, the largest coefficient of the
physicists' Hermite polynomial $H_n$ with signs retained: with
$k_n=\lfloor(2n+6-\lfloor\sqrt{8n+21}\rfloor)/8\rfloor$ and $d_n=n-4k_n$, the
exponent $d_n$ is the unique maximizer and $a(n)=n!\,2^{d_n}/((2k_n)!\,d_n!)$.
The distinction from the largest coefficient in absolute value is kept explicit
throughout.}

\entry{Asymptotics of OEIS A277364}
{generating-functions-and-asymptotics/oeis-sequence-asymptotics/a277364-bell-asymptotics}
{A277364\_asymptotics.pdf}{A277364\_asymptotics.zip}
{For $a(n)=\sum_{k\le n/2}S(n,k)$, partitions of an $n$-set into at most
$\lfloor n/2\rfloor$ blocks, proves $a(n)\sim B_n$ with an explicit first
correction: writing $r=W_0(n)$,
$a(n)=\exp\{n(r-1+1/r)-1\}(1+r)^{-1/2}\,(1-P(r)/n+O(r^2/n^2))$ for an explicit
rational $P$.  Also gives a parity-split equivalent for the omitted tail
$B_n-a(n)$.}

\entry{Term counts in the derivatives of power towers}
{generating-functions-and-asymptotics/oeis-sequence-asymptotics/power-tower-derivative-term-counts}
{article.pdf}{A293239\_research\_report.zip, A290268\_partial\_results.zip,
A281434\_exact\_algorithm\_and\_cubic\_growth.zip}
{One article on three OEIS sequences that ask the same question of three
functions: how many distinct monomials are in the fully collected $n$th
derivative of $x^x$ (A293239), $x^{x^2}$ (A290268) and $x^{x^x}$ (A281434)?
Part~I proves once the apparatus all three had established separately --- the
canonical differential-polynomial form $f^{(n)}=fP_n$ or $fx^{-n}P_n$ with its
first-order integer recurrence, the independence of logarithmic monomials that
makes the count an invariant of $f$ rather than of a differentiation algorithm,
the identity $a(n)=|\operatorname{supp}P_n|$, and the envelope bound.  Two
normalizations are needed, because neither clause covers all three in the frames
the parts use.  \textbf{None of the three headline questions is fully settled,
and the article says so in its abstract}: for $x^x$ the count
reduces exactly to zeros of the first-kind Lehmer--Comtet triangle A008296 and
$a(n)=\Theta(n^2)$, but the closed formula needs an isolated nonvanishing step
--- and separately the recurrence printed in the entry for $n>6$ is
\emph{disproved} at $n=8$; for $x^{x^2}$ the conjectured expression is proved to
be an upper bound and $a(n)=\Theta(n^2)$, but not the constant $\tfrac12$; for
$x^{x^x}$ explicit bounds give $a(n)=\Theta(n^3)$ but not $a(n)\sim\tfrac23n^3$,
and there the obstruction is not a nonvanishing statement but a density one,
since further zeros provably exist.  A closing part sets the three side by side
and records what they share beyond the framework: every cancellation any of them
can prove is a parity zero of the same falling factorial.  This is a thematic consolidation,
not a deduplication: the three were checked and are not duplicates, and all
fifty numbered results of the three originals are preserved unchanged.}

\entry{Iterated sums and products from one}
{generating-functions-and-asymptotics/oeis-sequence-asymptotics/a352969-correctness-and-growth}
{article.pdf}{A352969\_correctness\_and\_growth.zip}
{For $S_0=\{1\}$, $S_{n+1}=(S_n+S_n)\cup(S_n\cdot S_n)$ and $a_n=|S_n|$ (OEIS
A352969), proves correctness of the listed terms and the doubly exponential law
$2^{2^{\,n-r_n-4}}\le a_n\le 2^{2^{\,n-1}}$, hence
$\log\log a_n=n+O(\sqrt{\log n})$.  The step is simultaneous, not repeated
closure --- a distinction the report is careful about.}

\entry{Distinct third derivatives of a power tower}
{generating-functions-and-asymptotics/oeis-sequence-asymptotics/a199085-exponentiation-derivative-values}
{article.pdf}{A199085\_generating\_function.zip}
{For OEIS A199085, the number of distinct values taken by the third derivative
at $x=1$ over all fully parenthesized exponentiations of $n$ copies of $x$,
proves $a(1)=a(2)=1$ and $a(n)=\lfloor(n^2-2)/3\rfloor$ for $n\ge3$, with the
rational generating function $z/(1-z)+z^3/(1-z)^3-z^7/((1-z)^2(1-z^3))$.  The
proof is for all $n$, not an extrapolation: every attainable pair of second and
third derivatives is classified and a witness constructed for each admissible
pair.  Despite the shared vocabulary this is a different question from
\emph{power-tower-derivative-term-counts}, which fixes the function and varies
the derivative order; here the order is fixed at three and the parenthesization
varies.}

\subsection{Asymptotics, records and discrepancy}

\entry{The Exponential Scale of Gregory Sign Thresholds}
{generating-functions-and-asymptotics/gregory-coefficient-sign-thresholds}
{gregory\_sign\_threshold.pdf}{gregory\_sign\_threshold\_research.zip}
{For $(x/\log(1+x))^m=\sum_n G^{(m)}_nx^n$ and $\mathrm{Gr}(m)$ the least $N$
beyond which $(-1)^{n-1}G^{(m)}_n>0$, proves
$\log\mathrm{Gr}(m)=m+1-\gamma-(\zeta(2)+\zeta(3))/m+C_2/m^2+O(m^{-3})$ with
$C_2$ explicit, plus a finite algorithm for every coefficient of the complete
expansion; in particular $\mathrm{Gr}(m)/e^m\to e^{1-\gamma}$ and
$\mathrm{Gr}(m)/3^m\to0$.  This replaces Conjecture~6.1 of Xu and Zhao
(arXiv:2609.11072v1), whose printed lower bound already conflicts with its own
reported $\mathrm{Gr}(5)=113$; that value is independently certified here.}

\entry{Exact Records in the Square-Root-of-Three Walk}
{generating-functions-and-asymptotics/sqrt-three-walk-records}
{article.pdf}{sqrt3\_record\_recurrence.zip}
{Proves the four-phase record recurrence for
$S(N)=\sum_{n\le N}(-1)^{\lfloor n\sqrt3\rfloor}$ in the corrected form on
page~7 of Bruin and Fokkink, \emph{On the records and zeros of a deterministic
random walk} (arXiv:2503.11734v2), where it is experimental and attributed to
Arias de Reyna and van de Lune.  Includes minimality of every claimed first
passage, a rational generating function, closed forms, sharp logarithmic
extrema, and an extension to every even integer trace $D\ge4$.}

\entry{Sharp discrepancy for a family of binary substitutions}
{generating-functions-and-asymptotics/binary-substitution-discrepancy}
{article.pdf}{oeis\_discrepancy\_counterexamples.zip}
{For the fixed word of $0\mapsto1$, $1\mapsto101010$, the position sequences
A284365 and A284366 both carry a conjectured upper position error below~2, and
both conjectures are false: the first counterexamples are at ranks
$2{,}977{,}771$ and $8{,}933{,}313$, with common error
$(2977771\sqrt{21}-13645857)/2\approx2.00487$ and sharp common upper bound
$(6+\sqrt{21})/5\approx2.11652$.  Proves complete discrepancy limit intervals
for every substitution $0\mapsto1$, $1\mapsto(10)^m$, with infinite
counterexample families and exact recurrences for their indices.}

\entry{Exact eventual maximizers of Rudin--Shapiro autocorrelations}
{generating-functions-and-asymptotics/rudin-shapiro-autocorrelation-maximizers}
{article.pdf}{rudin\_shapiro\_maximizers.zip}
{For the aperiodic autocorrelation $C_m(k)$ of the level-$m$ Rudin--Shapiro
sequence, classifies the maximizing shift exactly for all large $m$:
$k_m^\ast=(2^{m+1}+(-1)^m)/3$, with $A_m=b_{m-2}$, settling the discrete
uniqueness-and-location conjecture of Tarnu (J.~Approx.\ Theory 287 (2023)
105866), repeated as Conjecture~1.5 of Choi--Tarnu (2025), including the limit
$k_m^\ast/2^m\to2/3$.  The certificate is a rational polytope argument with
exact arithmetic.}

\entry{Apéry arrays and conjectural zeta accelerations}
{generating-functions-and-asymptotics/apery-array-zeta-accelerations}
{article.pdf}{Apery\_OEIS\_Conjectures\_Proofs.zip, apery\_oeis\_conjectures.zip}
{Proves the entire shifted-diagonal $\zeta(3)$ family recorded as conjectural in
OEIS A143007, and both the super- and subdiagonal $\zeta(2)$ families of
A108625, at every nonnegative offset: for all $k\ge0$,
$\zeta(3)=H_k^{(3)}+\sum_{n\ge1}(2n+k)(3n^2+3nk+k^2)/
(n^3(n+k)^3T_{n-1,n+k-1}T_{n,n+k})$.  The main diagonals are the classical
Apéry sequences A005259 and A005258.  Proofs are finite telescoping identities
with rational lattice potentials and explicit limit estimates; no computer test
or unproved conjecture enters.  Adds rational error bounds, numerator and
denominator recurrences, a three-term recurrence along every shifted diagonal,
and sharp fixed-offset asymptotic constants.  Merged from two reports whose
boxed identities were character-identical up to macro names.}

\end{document}
